\chapter{Водород в движении}\label{ch:hydrogen}

\section{Диффузия}

Растворённый водород в цирконии очень подвижен. Он подчиняется закону Фика
\begin{equation}
\frac{\partial c}{\partial t}=D\,\nabla^2c,\qquad D=D_0\,e^{-Q_D/RT},
\end{equation}
с $D_0=7.9\cdot10^{-7}$~м$^2$/с и $Q_D=44.4$~кДж/моль. Характерное расстояние, на которое
водород уходит за время $t$, --- $\sqrt{2Dt}$ (рис.~\ref{fig:F17_Diffusion}).

\fig[0.9]{F17_Diffusion}{Диффузионная длина водорода за 10~с, минуту и 10~минут. За минуту при
300\,$^\circ$C водород проходит $\sim90$~мкм --- треть поля автомата.}

\begin{idea}
За время, пока металл остывает на градус (20~с при 3\,$^\circ$C/мин), водород пробегает десятки
микрон. Поэтому зёрна соревнуются за водород не каждое у себя, а все сразу: первое семейство,
начавшее выпадать, обедняет раствор во всём поле, и остальным уже нечего пересыщать. Это делает
порог резким (гл.~\ref{ch:nucleation}).
\end{idea}

На сетке диффузия считается через Фурье, как и упругость: за шаг $\delta t$ каждая гармоника
затухает точно,
\begin{equation}
\hat c(\mathbf k)\leftarrow\hat c(\mathbf k)\,e^{-Dk^2\delta t},
\end{equation}
что устойчиво при любом шаге.

\begin{exercise}
Посчитайте $D$ при 300\,$^\circ$C и при 200\,$^\circ$C и диффузионную длину за минуту.
Сравните с размером зерна (2.5--4.5~мкм) и с шагом между окружными пакетами у Lepine (57~мкм).
\end{exercise}

\section{Баланс водорода}

Пластинка забирает водород: в единице её площади (в двумерном сечении) сидит
$C_{\text{гидр}}\approx15\,660$~ppm металла (гл.~\ref{ch:problem}). Когда зародыш появляется в
клетке, модель проверяет, есть ли в окрестности радиусом $\sqrt{2D t_{\text{рост}}}$
($t_{\text{рост}}=10$~с) столько водорода сверх TSSD, чтобы вырастить пластинку минимальной
длины $L_{\min}=1.2$~мкм. Если нет --- зародыш не рождается. Водород пластинки снимается
сразу из её клеток, дальше яма заполняется диффузией.

\begin{exercise}
Пластинка $5\times0.6$~мкм. Сколько ppm$\cdot$мкм$^2$ водорода она содержит? Из какой площади
её надо собрать, если пересыщение над TSSD --- 20~ppm? Какой это радиус круга?
\end{exercise}

\section{Рост во времени}

\fig[0.82]{F18_Growth}{Рост пластинки с концов: водород стекается к пластинке (зелёные стрелки),
вокруг неё раствор обеднён; кончики (красные стрелки) удлиняются со скоростью, пропорциональной
пересыщению.}

Зародыш появляется длиной $L_{\min}$ и растёт с концов. Скорость кончика ограничена подводом
водорода: при кончике радиусом $h/2$ и пересыщении $c-c_{\TSSD}$ предельная скорость по диффузии
$\sim D/h\cdot(c-c_{\TSSD})/C_{\text{гидр}}$. Реальный рост медленнее (как в модели HNGD, Lacroix и др.),
поэтому
\begin{keyf}[Скорость кончика]
\[
v=k_{\text{tip}}\,\frac{D}{h}\,\frac{c-c_{\TSSD}}{C_{\text{гидр}}},\qquad k_{\text{tip}}=0.05\ \text{(допущение)} .
\]
\end{keyf}
Кончик останавливается у границы зерна, у соседней пластинки или при длине $L_{\max}=6$~мкм.
Упёршись в границу, он может продолжиться в соседнем зерне (или на соседней грани), если след
базисной плоскости там отличается не больше чем на $15^\circ$ --- так гидриды проходят сквозь
границы (Fang 2017, Son 2026).

\begin{exercise}
Скорость кончика при 300\,$^\circ$C и пересыщении 20~ppm --- в мкм/мин. За сколько минут
пластинка дорастёт от 1.2 до 5~мкм? Что это значит для сравнения скоростей охлаждения
0.3 и 15\,$^\circ$C/мин?
\end{exercise}

\section{Водород и напряжения}

Водород расширяет решётку (парциальный мольный объём $V_H=1.7\cdot10^{-6}$~м$^3$/моль), поэтому его
потенциал зависит от гидростатического напряжения $\sigma_h=\mathrm{tr}\,\sigma/3$:
\begin{equation}
\mu=\mu_0+RT\ln c-V_H\sigma_h\quad\Longrightarrow\quad c_{\text{равн}}\propto\exp\frac{V_H\sigma_h}{RT}.
\end{equation}
Водород стекает в растянутые места (эффект Горского); поток по Даттону и Пульсу
$J=-D\left[\nabla c-c\,\frac{V_H}{RT}\nabla\sigma_h\right]$ --- основа модели замедленного гидридного
растрескивания (DHC). В автомате это опция \code{stress\_diff}.

\begin{warnbox}[Проверено МКЭ]
Вокруг пластинки металл гидростатически \emph{сжат}: у кончика $\sigma_h$ на 540~МПа ниже
среднего. Водорода там на 17\% меньше --- кромка гидрида водород не стягивает, а отталкивает.
Нагрузка поднимает $\sigma_h$ у кромки радиальной пластинки на $0.4\,\sigma_\theta$; вместе с работой
следующей пластинки это даёт около $0.03\,^\circ$C/МПа --- втрое меньше измеренного. Поэтому
перенос водорода напряжением не объясняет ни порог, ни двухосность, и по умолчанию выключен.
\end{warnbox}

\begin{exercise}
Насколько вырастет равновесная концентрация водорода в зоне с $\sigma_h=+100$~МПа при 330\,$^\circ$C?
А при $\sigma_h=-540$~МПа?
\end{exercise}

\begin{codebox}
\textbf{Где в коде.} \code{ca\_kinetic.py}: \code{D0}, \code{QD}, \code{t\_grow}, баланс водорода у
зародыша (\code{avail}), рост кончиков \code{grow\_tips} (\code{grow\_kin}, \code{k\_tip}), переход через
границу \code{try\_cross} (\code{cross\_tol}), диффузия в конце шага (с \code{stress\_diff} --- через
переменные Слотбома).
\end{codebox}
